\begin{frame}{전체 트랜스포머 블록: 조각 조립}
  \begin{columns}[T]
    \begin{column}{0.56\textwidth}
      \begin{enumerate}
        \item \textbf{multi-head self-attention} (+ \textbf{잔차
          연결}, layer norm(Ba et al., 2016))
        \item \textbf{position-wise feed-forward network}: 각 위치의
          벡터를 \emph{독립적으로} 2층 MLP
          ($d_{model} \to 4\,d_{model} \to d_{model}$)에 통과
          (+ 잔차 연결, layer norm)
      \end{enumerate}
    \end{column}
    \begin{column}{0.4\textwidth}
      \centering
      \includegraphics[height=4.6cm]{ch12_3_transformer_block.png}\\[0.2em]
      {\footnotesize 인코더/디코더 블록 구조}
    \end{column}
  \end{columns}
\note{\begin{itemize}
\item 12.1$\sim$12.3의 조각(Q/K/V 투영, scaled dot-product,
  self/multi-head, PE, causal mask)을 조립하면
  \kb{트랜스포머 블록(transformer block)}이 완성:
\item \kb{잔차 연결 $y = x + F(x)$}는 ResNet(He et al.,
      2015)의 ``단축로''과 \textbf{같은 구조} — 6$\sim$12층을
      쌓아 그래디언트 경로가 길어지는 상황에서 중요.
\item \textbf{layer norm}은 \textbf{batch norm}(Ioffe \&
      Szegedy, 2015)과 비슷한 ``층마다 스케일 정규화'' 역할
      (Week 6.2 정규화 논리 참조).
\item {\footnotesize 인코더 블록 = MHA + FFN(각각 Add\&Norm),
  6회 쌓기. 디코더 블록 = masked self-attention +
  cross-attention(인코더 출력) + FFN, 6회 쌓기. 입력/출력
  임베딩에 PE를 더하고, 디코더 출력은 linear + softmax로
  \textbf{토큰 확률 분포}를 낸다.}
\end{itemize}}
\end{frame}
